% fig_rebuild_02_onehot_qubo_encoding.tex
% 候选图 2 / one-hot QUBO 编码示意（Q1/Q2 用 n=15，本图用 n=5 演示更清晰）
% 编译：xelatex fig_rebuild_02_onehot_qubo_encoding.tex
\documentclass[border=10pt]{standalone}
\usepackage{ctex}
\usepackage{amsmath,amssymb}
\usepackage{tikz}
\usetikzlibrary{positioning,arrows.meta,matrix,decorations.pathreplacing,calc}

\begin{document}
\begin{tikzpicture}[
    font=\small,
    cell/.style={rectangle, draw=gray!60, minimum size=8mm, align=center},
    one/.style={fill=blue!18},
    rowsum/.style={rectangle, draw=orange!80, fill=orange!10, minimum size=8mm, align=center},
    colsum/.style={rectangle, draw=green!55!black, fill=green!12, minimum size=8mm, align=center},
    arrow/.style={-{Latex[length=2mm]}, thick, gray!70!black},
]

% ====================== 左：决策矩阵 X = (x_{i,p}) ======================
\node[align=center, font=\bfseries\small]
  (title_left) at (0, 4.6) {(a) one-hot 决策矩阵 $X = (x_{i,p})$};

% 列标签：位置 p
\foreach \p [count=\j from 1] in {1,2,3,4,5}
  \node[align=center, font=\small]
    at ({(\j-0.5)*0.9}, 3.5) {$p\!=\!\p$};

% 行标签：客户 i
\foreach \i [count=\k from 1] in {1,2,3,4,5}
  \node[align=center, font=\small]
    at (-0.7, {3.0 - (\k-0.5)*0.9}) {$i\!=\!\i$};

% 矩阵单元（演示一条合法路线 1→3→5→2→4）
\foreach \i [count=\ii from 1] in {1,2,3,4,5} {
  \foreach \p [count=\pp from 1] in {1,2,3,4,5} {
    \pgfmathtruncatemacro{\onepos}{0}
    \pgfmathtruncatemacro{\onepos}{
      ifthenelse((\ii==1 && \pp==1) || (\ii==3 && \pp==2) ||
                 (\ii==5 && \pp==3) || (\ii==2 && \pp==4) ||
                 (\ii==4 && \pp==5), 1, 0)}
    \ifnum\onepos=1
      \node[cell, one]   at ({(\pp-0.5)*0.9}, {3.0 - (\ii-0.5)*0.9}) {1};
    \else
      \node[cell]        at ({(\pp-0.5)*0.9}, {3.0 - (\ii-0.5)*0.9}) {0};
    \fi
  }
}

% 行约束（每行 sum = 1）
\foreach \i in {1,...,5}
  \node[rowsum] at (5.0, {3.0 - (\i-0.5)*0.9}) {$=\!1$};
\node[font=\footnotesize, color=orange!90!black, align=left]
  at (6.5, 1.35)
  {\textbf{行约束}\\每个客户\\恰好访问\\1 次};

% 列约束（每列 sum = 1）—— 放到矩阵下方留出空隙
\foreach \p [count=\pp from 1] in {1,2,3,4,5}
  \node[colsum] at ({(\pp-0.5)*0.9}, -2.1) {$=\!1$};
\node[font=\footnotesize, color=green!50!black, align=center]
  at (2.0, -3.0) {\textbf{列约束}：每个位置恰好分配 1 个客户};

% ====================== 右：QUBO 能量与解码 ======================
\begin{scope}[xshift=110mm]
\node[align=center, font=\bfseries\small]
  at (0, 4.6) {(b) QUBO 能量分解 \& 解码};

\node[align=left, font=\small]
  at (0, 3.2)
  {$E(\mathbf x) \;=\; A\!\cdot\! H_{\text{const}}(\mathbf x)\;+\; H_{\text{dist}}(\mathbf x)$};

\node[align=left, font=\footnotesize, text width=80mm]
  at (0, 2.0)
  {$H_{\text{const}}=\displaystyle\sum_{i=1}^{n}\!\Big(\sum_{p}x_{ip}-1\Big)^{\!2}\!+\!\sum_{p=1}^{n}\!\Big(\sum_{i}x_{ip}-1\Big)^{\!2}$};

\node[align=left, font=\footnotesize, text width=80mm]
  at (0, 0.5)
  {$H_{\text{dist}}=\displaystyle\sum_{p=1}^{n-1}\sum_{i\ne j}T_{ij}\,x_{i,p}\,x_{j,p+1}\;+\;\text{depot}$ 项};

% 解码箭头
\draw[arrow] (0, -0.5) -- (0, -1.3);

\node[align=center, font=\small,
      rectangle, draw=blue!75, fill=blue!8, rounded corners=2pt,
      minimum width=70mm, minimum height=10mm]
  at (0, -1.9)
  {\textbf{解码：}$x_{i,p}=1 \Rightarrow$ 第 $p$ 位访问客户 $i$};

\node[align=center, font=\small\ttfamily,
      rectangle, draw=green!55!black, fill=green!10, rounded corners=2pt,
      minimum width=70mm, minimum height=8mm]
  at (0, -3.2) {路线: 0 $\to$ 1 $\to$ 3 $\to$ 5 $\to$ 2 $\to$ 4 $\to$ 0};
\end{scope}

% ====================== 底部说明 ======================
\node[align=center, font=\footnotesize, text width=170mm]
  at (8.0, -5.0)
  {\textbf{比特预算公式：}$n_{\text{var}}=n^{2}$（本文图示 $n\!=\!5$ 仅作演示；
   Q1/Q2 实际 $n\!=\!15 \Rightarrow 225$ 比特，全部 $\le 550$ 真机比特上限）};

\end{tikzpicture}
\end{document}
